\chapter{Etude Prealable}
\label{chap:etude-prealable}

\section{Introduction}
Ce chapitre presente le cadre general du projet de fin d'etudes realise au sein de l'entreprise Tritux. Il introduit le contexte, la problematique, les objectifs, la methode de travail et les acteurs impliques dans la solution.

\section{Presentation de l'organisme d'accueil}
\subsection{Presentation de Tritux}
Tritux est une entreprise active dans le developpement de solutions numeriques et l'accompagnement des organisations dans leur transformation digitale. Son orientation est basee sur l'innovation, l'automatisation des processus et l'amelioration de la productivite des equipes metier.

\subsection{Contexte du projet chez Tritux}
Dans le cadre de ses besoins internes et de ses activites RH, Tritux souhaite disposer d'une plateforme qui simplifie la gestion des CV, ameliore leur qualite de presentation et automatise les actions repetitives de traitement des candidatures.

\section{Contexte et problematique}
Le recrutement moderne repose sur des volumes importants de candidatures et sur des exigences elevees en termes de qualite de presentation des profils. Dans un traitement manuel classique, plusieurs limites apparaissent :
\begin{itemize}[label=--]
    \item surcharge de traitement due a la lecture manuelle des CV,
    \item heterogeneite des formats et perte de temps dans la mise en forme,
    \item difficulte d'adapter rapidement les CV a une charte de branding,
    \item absence d'un mecanisme unifie de notification des candidats.
\end{itemize}

Ces contraintes rallongent les delais de selection, augmentent le risque d'erreurs et reduisent l'efficacite des equipes RH. D'ou la necessite de concevoir une plateforme intelligente capable d'automatiser l'analyse, l'optimisation et la transformation des CV.

\section{Objectifs du projet}
\subsection{Objectif general}
Concevoir et developper une plateforme intelligente de branding et d'optimisation de CV avec notification automatisee des candidats.

\subsection{Objectifs specifiques}
\begin{itemize}[label=--]
    \item importer des CV depuis une interface centralisee,
    \item extraire automatiquement les informations principales du CV,
    \item optimiser et restructurer le contenu via des mecanismes d'intelligence artificielle,
    \item transformer le CV selon des templates personnalises,
    \item offrir une previsualisation avant generation finale en DOCX,
    \item notifier automatiquement les candidats selectionnes par email.
\end{itemize}

\section{Etude de l'existant}
\subsection{Processus actuel}
Le flux de traitement classique des CV s'appuie principalement sur des operations manuelles : reception des candidatures, lecture individuelle, correction de forme, adaptation visuelle et communication separee avec les candidats.

\subsection{Limites de l'existant}
\begin{itemize}[label=--]
    \item absence d'une centralisation complete des operations,
    \item faible standardisation de la qualite visuelle des CV,
    \item manque d'automatisation pour les taches repetitives,
    \item suivi difficile des etats de transformation et de notification.
\end{itemize}

\section{Solution proposee}
La solution proposee est une plateforme web composee d'un backend de traitement et d'une interface utilisateur moderne. Elle prend en charge l'import des CV, l'extraction d'informations, l'optimisation du contenu, l'application de templates et l'envoi de notifications.

La plateforme cible principalement deux profils d'utilisateurs :
\begin{itemize}[label=--]
    \item \textbf{Administrateur RH} : pilote l'ensemble du flux de transformation et de suivi,
    \item \textbf{Gestionnaire Recrutement} : exploite les resultats et declenche les notifications.
\end{itemize}

\section{Diagramme de cas d'utilisation global}
La figure suivante presente une vue globale des interactions entre les acteurs et les principales fonctionnalites de la plateforme.

\begin{figure}[H]
\centering
\begin{tikzpicture}[>=latex, node distance=1.2cm]
    % System boundary
    \node[draw, minimum width=12cm, minimum height=8cm] (system) {};
    \node[above left] at (system.north east) {Plateforme intelligente CV};

    % Actors
    \node (rh) at (-7,2.5) {\textbf{Administrateur RH}};
    \draw (-7,1.9) circle (0.3);
    \draw (-7,1.6) -- (-7,0.8);
    \draw (-7,1.35) -- (-7.5,1.1);
    \draw (-7,1.35) -- (-6.5,1.1);
    \draw (-7,0.8) -- (-7.4,0.2);
    \draw (-7,0.8) -- (-6.6,0.2);

    \node (gr) at (-7,-1.2) {\textbf{Gestionnaire Recrutement}};
    \draw (-7,-1.8) circle (0.3);
    \draw (-7,-2.1) -- (-7,-2.9);
    \draw (-7,-2.35) -- (-7.5,-2.6);
    \draw (-7,-2.35) -- (-6.5,-2.6);
    \draw (-7,-2.9) -- (-7.4,-3.5);
    \draw (-7,-2.9) -- (-6.6,-3.5);

    % Use cases
    \node[draw, ellipse, minimum width=4.3cm] (u1) at (-0.5,2.6) {Importer des CV};
    \node[draw, ellipse, minimum width=4.3cm] (u2) at (3.2,1.4) {Extraire les informations};
    \node[draw, ellipse, minimum width=4.8cm] (u3) at (0.7,0.2) {Optimiser le contenu du CV};
    \node[draw, ellipse, minimum width=5.2cm] (u4) at (2.2,-1.1) {Adapter le CV a un template};
    \node[draw, ellipse, minimum width=4.5cm] (u5) at (0.2,-2.4) {Previsualiser le resultat};
    \node[draw, ellipse, minimum width=5.3cm] (u6) at (2.8,-3.8) {Generer le CV final (DOCX)};
    \node[draw, ellipse, minimum width=5.3cm] (u7) at (0.2,-5.1) {Notifier les candidats par email};

    % Associations
    \draw[-] (-6.7,1.1) -- (u1.west);
    \draw[-] (-6.7,1.1) -- (u2.west);
    \draw[-] (-6.7,1.1) -- (u3.west);
    \draw[-] (-6.7,-2.6) -- (u4.west);
    \draw[-] (-6.7,-2.6) -- (u5.west);
    \draw[-] (-6.7,-2.6) -- (u6.west);
    \draw[-] (-6.7,-2.6) -- (u7.west);
\end{tikzpicture}
\caption{Diagramme de cas d'utilisation global de la plateforme}
\label{fig:usecase-global}
\end{figure}

\section{Methodologie adoptee}
Une approche iterative et incrementale a ete retenue afin de livrer progressivement les fonctionnalites, integrer les retours des encadrantes et reduire les risques techniques.

Les grandes etapes du travail sont :
\begin{enumerate}
    \item etude de l'existant et specification des besoins,
    \item conception fonctionnelle et technique,
    \item developpement des modules de la plateforme,
    \item tests, integration et validation.
\end{enumerate}

\section{Plan du rapport}
Le rapport est organise comme suit :
\begin{itemize}[label=--]
    \item \textbf{Chapitre 1} : etude prealable du projet,
    \item \textbf{Chapitre 2} : analyse des besoins et specification,
    \item \textbf{Chapitre 3} : conception de la solution,
    \item \textbf{Chapitre 4} : realisation et validation,
    \item \textbf{Conclusion generale} : bilan et perspectives.
\end{itemize}

\section{Conclusion}
Ce chapitre a presente le contexte de realisation du projet, les limites de l'existant, les objectifs cibles et la vision globale de la solution. Le chapitre suivant detaille l'analyse des besoins fonctionnels et non fonctionnels.
